\exnumber{02}
\extitle{UX VI - Ex 3 - Animation des wireframes}
\exfiles{ex02/fichier Figma, nomme Name\_FirstName\_DeliverableName\_Date\_VersionNumber}
\exauthorize{Those named in the exercise}
\makeheaderfiles

\textbf{Instructions :}

Animate your wireframes on Figma.

Think about multistate!!

Cet exercice se fait seul.

\textbf{Livrables :}

\begin{itemize}
\item An animated prototype on Figma
\item Déposez le fichier Figma sur Github, nommé ainsi : « Name\_FirstName\_DeliverableName\_Date\_VersionNumber ». Exemple : « Doe\_John\_MidFiPrototype\_05242024\_V1 ».
\item Ajoutez un titre dans le fichier.
\item Soignez l'aspect global du fichier. Il doit être clair, simple et lisible. Vous pouvez regarder canva.com, sans surcharger le dessin de détails.
\end{itemize}

\textbf{Ressources :}

\begin{itemize}
\item Efficiently Manage Your Designs - How Mockplus State Page Helps To View \& Manage Multiple State \url{https://help.mockplus.com/p/372}
\end{itemize}

